\documentclass[12pt]{article}
\usepackage{amsfonts,amsthm,amsmath,amssymb,mathtools}
\usepackage{mathrsfs}
\usepackage{enumitem}
\usepackage{tikz-cd}
\usepackage{geometry}
\usepackage{graphicx}
\usepackage{mlmodern}

\geometry{letterpaper, portrait, margin=1in}

\setcounter{section}{0}

\renewcommand{\thesection}{\S\arabic{section}}
\renewcommand{\thesubsection}{\arabic{section}}
\DeclareMathOperator{\im}{im}
\DeclareMathOperator{\id}{id}
\DeclareMathOperator{\ad}{ad}
\DeclareMathOperator{\tr}{tr}
\newcommand{\gl}{\mathfrak{gl}}
\newcommand{\so}{\mathfrak{so}}
\renewcommand{\sl}{\mathfrak{sl}}
\renewcommand{\th}{^\mathrm{th}}

\newtheorem{thm}{Theorem}
\newenvironment{theorem}[1]{
	\renewcommand{\thethm}{#1}
	\begin{thm}
		}{\end{thm}}

\newtheorem{lma}{Lemma}
\newenvironment{lemma}[1]{
	\renewcommand{\thelma}{#1}
	\begin{lma}
		}{\end{lma}}

\theoremstyle{definition}
\newtheorem{ex}{}
\newenvironment{exercise}[1]{
	\renewcommand{\theex}{#1}
	\begin{ex}
		}{\end{ex}}

\title{Lie Theory: Chapter 5}
\author{Connor Mooneyhan}
\date{}

\begin{document}

\maketitle

\begin{ex}
	Let $V$ be a vector space, $F$ be a field, and $A$ be a subalgebra of $\gl(V)$.
	Give an example of a function $f : A \to F$ whose eigenspace is zero.
\end{ex}
\begin{proof}
	Define $f$ to be the map sending every element to 1.
	We know that the zero map $0 \in A$ since $A$ is a subalgebra, so this would mean that if $a(v) = f(a)v$ for all $a \in A$, in particular we have $0 = 0(v) = f(0)v = 1v = v$, so the weight space/eigenspace $V_f$ is zero.
\end{proof}

\begin{ex}
	Let \begin{equation*}
		A = \left\lbrace \begin{pmatrix} a & b \\ 0 & a + b \end{pmatrix} : a,b \in \mathbb{C} \right\rbrace.
	\end{equation*}
	\begin{enumerate}[label=(\alph*)]
		\item Prove that $A$ is an abelian subalgebra of $\gl(2,\mathbb{C})$.
		\item Let $\lambda : A \to \mathbb{C}$ be defined by $\lambda \begin{pmatrix}a & b \\ 0 & a+b\end{pmatrix} = a + b.$
		      Find the weight space $V_\lambda$ associated to $\lambda$.
	\end{enumerate}
\end{ex}
\begin{proof}\
	\begin{enumerate}[label=(\alph*)]
		\item \begin{align*}
			          & \left[\begin{pmatrix}a & b \\ 0 & a+b\end{pmatrix}, \begin{pmatrix}c & d \\ 0 & c+d\end{pmatrix}\right]                                                                             \\
			      =\  & \begin{pmatrix}a & b \\ 0 & a+b\end{pmatrix}\begin{pmatrix}c & d \\ 0 & c+d\end{pmatrix} - \begin{pmatrix}c & d \\ 0 & c+d\end{pmatrix}\begin{pmatrix}a & b \\ 0 & a+b\end{pmatrix} \\
			      =\  & \begin{pmatrix}ac & ad + bc + bd \\ 0 & ac + ad + bc + bd\end{pmatrix} - \begin{pmatrix}ac & bc + ad + bd \\ 0 & ac + bc + ad + bd\end{pmatrix}                                     \\
			      =\  & \begin{pmatrix}0 & 0 \\ 0 & 0\end{pmatrix}
		      \end{align*}
		\item Suppose that $v \in V$ and $av = \lambda(a)v$ for all $a \in A$.
		      Then writing $a = \begin{pmatrix}a & b \\ 0 & a + b\end{pmatrix}$, \begin{align*}
			      \begin{pmatrix}av_1 + bv_1 \\ av_2 + bv_2\end{pmatrix} & = (a + b)\begin{pmatrix}v_1\\v_2\end{pmatrix}                                        \\
			                                                             & = (a + b)v                                                                           \\
			                                                             & = \lambda(a)v                                                                        \\
			                                                             & = av                                                                                 \\
			                                                             & = \begin{pmatrix}a & b \\ 0 & a + b\end{pmatrix}\begin{pmatrix}v_1\\v_2\end{pmatrix} \\
			                                                             & = \begin{pmatrix}av_1 + bv_2 \\ av_2 + bv_2\end{pmatrix}.
		      \end{align*}
		      This leaves us with two equations: $av_1 + bv_1 = av_1 + bv_2$ and $av_2 + bv_2 = av_2 + bv_2$.
		      Flearly the second equation is trivial, but the first tells us that $bv_1 = bv_2$.
		      This equality holds if $b = 0$, but we need the equation to hold for all possible values of $b$, so we are left with $v_1 = v_2$.
		      Thus \begin{equation*}
			      V_\lambda = \left\lbrace \begin{pmatrix}x\\x\end{pmatrix} : x \in \mathbb{C} \right\rbrace.
		      \end{equation*}
	\end{enumerate}
\end{proof}

\begin{ex}\
	\begin{enumerate}[label=(\alph*)]
		\item Let $L$ be a Lie subalgebra of $\gl(V)$.
		      Suppose that there is a basis of $V$ such that, with respect to this basis, every $x \in L$ is represented by a strictly upper triangular matrix.
		      Show that $L$ is isomorphic to a Lie subalgebra of $n(n, F)$ and hence that $L$ is nilpotent.
		\item Prove a result analagous to (i) in the case where there is a basis of $V$ such that with respect to this basis every $x \in L$ is represented by an upper triangular matrix.
	\end{enumerate}
\end{ex}
\begin{proof}\
	\begin{enumerate}[label=(\alph*)]
		\item Let $f : L \to n(n,F)$ be defined by taking each $x \in L$ to its strictly upper triangular matrix under the aforementioned basis.
		      $\im f$ is closed under linear combinations since $ax + by \in L$ by virtue of $L$ being a Lie subalgebra, and thus $f(ax + by) = af(x) + bf(y) \in n(n,F) \in \im f$ (which also shows that $f$ is a linear map), and $\im f$ is closed under brackets by similar logic, and likewise with showing that $f$ carries brackets to brackets, showing it is a homomorphism.
		      Of course, $f$ is automatically surjective onto its image, so the only thing left to show is that it is injective.
		      Indeed, if $f(x) = f(y)$, then the matrices for $x$ and $y$ are the same \textit{over the same basis}, so $x = y$, and $f$ is an isomorphism onto $\im f$, which is a subalgebra of $\gl(V)$.
		\item Nothing in the proof of (a) uses the fact that the matrices were \textit{strictly} upper triangular.
		      The proof is thus identical, modulo the word ``strictly''.
	\end{enumerate}
\end{proof}

\begin{ex}
	Let $L$ be a Lie algebra and let $A$ be a subalgebra of $L$.
	The \textit{normalizer} of $A$, denoted $N_L(A)$, is defined by \begin{equation*}
		N_L(A) = \lbrace x \in L : [x,a] \in A \text{ for all } a \in A\rbrace.
	\end{equation*}
	Prove that $N_L(A)$ is a subalgebra of $L$ and that $N_L(A)$ contains $A$.
	Show moreover that $N_L(A)$ is the largest subalgebra of $L$ in which $A$ is an ideal.
\end{ex}
\begin{proof}
	Certainly $A \subseteq N_L(A)$ since $A$ is a subalgebra and thus $[a, b] \in A$ whenever $a, b \in A$.
	If $x, y \in N_L(A)$ and $\alpha, \beta \in F$, then given $a \in A$, \begin{align*}
		[\alpha x + \beta y, a] & = \alpha[x, a] + \beta[y, a] \in A
	\end{align*}
	since both $[x, a] \in A$ and $[y, a] \in A$ and $A$ is a subalgebra, so $N_L(A)$ is closed under linear combinations.
	Furthermore, \begin{align*}
    [[x, y], a] & = \underbrace{[\overbrace{[x, a]}^{\in A}, y]}_{\in A} + \underbrace{[\overbrace{[a, y]}^{\in A}, x]}_{\in A} \hspace{2em}\text{(by Jacobi)}
	\end{align*}
  where every instance of ``$\in A$'' comes from taking the bracket of an element in $N_L(A)$ with an element of $A$.
  Thus $N_L(A)$ is a subalgebra of $L$.

  Now suppose $B$ is a subalgebra of $L$ which contains $N_L(A)$ as a proper subset, and let $b \in B \setminus N_L(A)$.
  Then there exists an $a \in A$ such that $[b, a] \notin A$ (since otherwise $b$ would be in $N_L(A)$), so $A$ is not an ideal of $B$.
  Thus $N_L(A)$ is the largest subalgebra in which $A$ is an ideal.
\end{proof}

\begin{ex}
  Let $L = \gl(n, \mathbb{C})$ and let $A$ be the subalgebra of $L$ consisting of all diagonal matrices.
  Show that $N_L(A) = A$.
\end{ex}

\begin{ex}
  Show that if $a, y \in \gl(V)$ are linear maps, then for any $m \geq 1$ \begin{equation*}
    ay^m = y^m a + \sum_{k=1}^m \begin{pmatrix}m \\ k\end{pmatrix} y^{m-k}a_k
  \end{equation*}
  where $a_1 = [a,y]$, $a_2 = [a_1, y] = [[a,y],y]$, and generally $a_k = [a_{k-1},y]$.
  Deduce that \begin{equation*}
    \ad y^m = \sum_{k=1}^m (-1)^{k+1}\begin{pmatrix}m \\ k\end{pmatrix}y^{m-k}(\ad y)^k.
  \end{equation*}
  In this formula, the power $y^{m-k}$ acts on $\gl(V)$ as the linear map which sends $x \in \gl(V)$ to $y^{m-k}x$.
\end{ex}

\begin{ex}
  This exercise gives an alternative approach to Proposition 5.7.
  Recall that $V$ is a complex vector space and $x,y \in \gl(V)$ are linear maps such that $[x,y]$ commutes with both $x$ and $y$.
  \begin{enumerate}[label=(\roman*)]
    \item Let $z = [x,y]$.
      Show that $\tr z^m = 0$ for all $m \geq 1$.
    \item Let $\lambda_1, \dots, \lambda_n$ be the eigenvalues of $[x, y]$.
      Show that $\lambda_i=0$ for all $i$, and deduce that $[x,y]$ is nilpotent.
      \textit{Hint:} use the Vandermonde determinant.
  \end{enumerate}
\end{ex}

\end{document}
